\documentclass[aspectratio=169,10pt]{beamer}
\makeatletter\def\input@path{{../../estilo/}{../estilo/}{estilo/}}\makeatother
\usepackage{devmobile}
\modulo[Flutter]{10}{Persistência Local}{shared\_preferences, armazenamento seguro, arquivos, sqflite e Drift com migrações e testes}

\begin{document}
\capa

\begin{frame}{Objetivos do módulo}
  \begin{conceito}[Ao final deste módulo você será capaz de]
    \begin{itemize}
      \item Escolher onde guardar cada tipo de dado no dispositivo.
      \item Salvar preferências com \cod{SharedPreferencesAsync} e segredos com \cod{flutter\_secure\_storage}.
      \item Ler e gravar arquivos com \cod{path\_provider}.
      \item Usar SQLite diretamente com \textbf{sqflite}.
      \item Modelar tabelas, consultas reativas e migrações com \textbf{Drift}.
      \item Tornar o banco local a fonte da verdade de um app \emph{offline-first}.
    \end{itemize}
  \end{conceito}
\end{frame}

\sumario

% =============================================================================
%  Fundamentos compartilhados --- Módulo 10: persistência, cache e sincronização
%  Incluído pelas trilhas Kotlin e Flutter com % =============================================================================
%  Fundamentos compartilhados --- Módulo 10: persistência, cache e sincronização
%  Incluído pelas trilhas Kotlin e Flutter com % =============================================================================
%  Fundamentos compartilhados --- Módulo 10: persistência, cache e sincronização
%  Incluído pelas trilhas Kotlin e Flutter com \input{../fundamentos/f10-persistencia}
%  (modelagem relacional, SQL, transações e índices estão no Módulo 09)
% =============================================================================
\section{Fundamentos: persistência, cache e sincronização}

\begin{frame}{Memória e armazenamento: velocidade \emph{versus} permanência}
  \begin{columns}[c]
    \begin{column}{0.5\textwidth}
      \centering
      \begin{tikzpicture}[every node/.style={font=\tiny}]
        \foreach \y/\w/\texto/\cor in {3/3.0/{registradores e cache: $\sim$1\,ns}/azulpetroleo,2.25/3.6/{memória RAM: $\sim$100\,ns}/azulpetroleo!80,1.5/4.2/{\emph{flash}: dezenas de µs}/ambar!85,0.75/4.8/{rede móvel: dezenas a centenas de ms}/cinzamedio!70}{
          \fill[\cor,rounded corners=2pt] (-\w/2,\y) rectangle (\w/2,\y+0.6);
          \node[text=white,font=\tiny\bfseries] at (0,\y+0.3) {\texto};
        }
        \draw[seta] (-2.75,3.55) -- node[rotulo,left,text width=1cm,align=right]{mais rápido e menor} (-2.75,0.8);
        \node[rotulo] at (0,0.45) {RAM: volátil \quad \emph{flash}: persistente};
\end{tikzpicture}
    \end{column}
    \begin{column}{0.48\textwidth}
      \begin{definicao}[Persistência]
        Guardar dados de forma que sobrevivam ao fim do processo, ao desligar do aparelho e à falta de rede. No celular, isso significa gravar na memória \emph{flash}.
      \end{definicao}
      \begin{conceito}[Por que persistir em um app]
        O sistema encerra apps em segundo plano (Módulo 07), a rede falha e o usuário espera abrir o app e ver seus dados na hora.
      \end{conceito}
      \begin{pordentro}[A memória \emph{flash}]
        Grava em blocos e se desgasta com o uso. Muitas gravações pequenas custam mais que uma grande: agrupe operações em transações (Módulo 09).
      \end{pordentro}
    \end{column}
  \end{columns}
\end{frame}

\begin{frame}{Onde o sistema guarda os arquivos do app}
  \scriptsize
  \begin{tabularx}{\textwidth}{@{}p{2.5cm} p{2.6cm} p{3.4cm} >{\tiny}L@{}}
    \toprule
    \cabecalho{Finalidade} & \cabecalho{Android} & \cabecalho{iOS} & \scriptsize\cabecalho{Flutter (\cod{path\_provider})} \\
    \midrule
    Dados internos do app & \cod{filesDir}, \cod{databases/} & \cod{Library/Application Support} & \cod{getApplicationSupportDirectory()} \\
    Documentos do usuário & \cod{filesDir} ou MediaStore & \cod{Documents} & \cod{getApplicationDocumentsDirectory()} \\
    Cache descartável & \cod{cacheDir} & \cod{Library/Caches} & \cod{getApplicationCacheDirectory()} \\
    Temporários & \cod{cacheDir} & \cod{tmp} & \cod{getTemporaryDirectory()} \\
    \bottomrule
  \end{tabularx}
  \normalsize
  \vspace{3pt}
  \begin{columns}[T]
    \begin{column}{0.32\textwidth}
      \begin{definicao}[\emph{Sandbox} de arquivos]
        Cada app tem pastas privadas que outros apps não leem. Desinstalar o app apaga tudo.
      \end{definicao}
    \end{column}
    \begin{column}{0.32\textwidth}
      \begin{pordentro}[Cópia de segurança]
        O Android copia até 25\,MB por app para a conta Google; o iOS inclui \cod{Documents} e \cod{Library}, exceto \cod{Caches}.
      \end{pordentro}
    \end{column}
    \begin{column}{0.32\textwidth}
      \begin{atencao}[Cache some]
        O sistema apaga caches quando falta espaço. Nunca guarde ali o que não puder baixar de novo.
      \end{atencao}
    \end{column}
  \end{columns}
\end{frame}

\begin{frame}{Do objeto à tabela}
  \begin{columns}[c]
    \begin{column}{0.46\textwidth}
      \centering
      \begin{tikzpicture}[every node/.style={font=\tiny}]
        \node[caixaescura,font=\tiny,align=left,text width=2.2cm] (o) at (0,0) {\textbf{objeto} \cod{Reserva}\\ \cod{bancada: Bancada}\\ \cod{leituras: List<...>}};
        \node[caixa,font=\tiny,minimum width=1.7cm] (t1) at (3.3,0.9) {tabela \cod{reserva}};
        \node[caixa,font=\tiny,minimum width=1.7cm] (t2) at (3.3,0) {tabela \cod{bancada}};
        \node[caixa,font=\tiny,minimum width=1.7cm] (t3) at (3.3,-0.9) {tabela \cod{leitura}};
        \draw[setaambar] (o.east) -- (t1.west); \draw[setaambar] (o.east) -- (t2.west); \draw[setaambar] (o.east) -- (t3.west);
        \node[rotulo,text width=4.6cm] at (1.6,-1.65) {referências viram chaves estrangeiras; listas viram outra tabela};
      \end{tikzpicture}
      \begin{definicao}[Descompasso objeto-relacional]
        Objetos têm referências, listas e herança; tabelas têm linhas, colunas e chaves. Traduzir de um mundo para o outro é o trabalho do mapeamento.
      \end{definicao}
    \end{column}
    \begin{column}{0.52\textwidth}
      \scriptsize
      \begin{tabularx}{\linewidth}{@{}l L@{}}
        \toprule
        \cabecalho{Abordagem} & \cabecalho{Como é} \\
        \midrule
        SQL à mão & você escreve o SQL e converte as linhas (\cod{sqflite}) \\
        SQL verificado & você escreve o SQL; a ferramenta confere na compilação e gera o código (Room, Drift) \\
        ORM completo & a biblioteca gera o SQL a partir dos objetos \\
        Banco de objetos & guarda objetos sem tabelas (ObjectBox, Realm) \\
        \bottomrule
      \end{tabularx}
      \normalsize
      \begin{dica}
        Room e Drift ficam no meio-termo: SQL explícito, mas com tipos e erros apontados antes de o app rodar.
      \end{dica}
      \begin{atencao}[O problema N+1]
        Buscar uma lista e depois uma consulta por item gera N+1 idas ao banco. Use \cod{JOIN} ou relações carregadas de uma vez.
      \end{atencao}
    \end{column}
  \end{columns}
\end{frame}

\begin{frame}{O esquema também tem versões}
  \centering
  \begin{tikzpicture}[node distance=16mm,every node/.style={font=\scriptsize}]
    \node[caixa,minimum width=1.4cm] (v1) {v1};
    \node[caixa,minimum width=1.4cm,right=of v1] (v2) {v2};
    \node[caixa,minimum width=1.4cm,right=of v2] (v3) {v3};
    \node[caixaescura,minimum width=1.4cm,right=of v3] (v4) {v4 (atual)};
    \draw[seta] (v1) -- node[rotulo,above]{+ coluna} (v2);
    \draw[seta] (v2) -- node[rotulo,above]{nova tabela} (v3);
    \draw[seta] (v3) -- node[rotulo,above]{renomeia} (v4);
    \draw[setaambar,rounded corners] (v1.south) -- ++(0,-4mm) -| node[rotulo,pos=0.25,below]{usuário que pulou versões: o app aplica 1$\rightarrow$2$\rightarrow$3$\rightarrow$4 em sequência} (v4.south);
  \end{tikzpicture}
  \raggedright
  \begin{columns}[T]
    \begin{column}{0.49\textwidth}
      \begin{definicao}[Migração]
        Programa que leva o banco de uma versão do esquema para a seguinte \textbf{sem perder os dados} do usuário. Cada versão publicada do app pode encontrar qualquer versão anterior do banco.
      \end{definicao}
    \end{column}
    \begin{column}{0.49\textwidth}
      \begin{conceito}[Regras de ouro]
        Migrações só andam para a frente e nunca são editadas depois de publicadas; teste a partir de \textbf{cada} versão antiga; apagar e recriar o banco só serve para dados de cache.
      \end{conceito}
    \end{column}
  \end{columns}
\end{frame}

\begin{frame}{Cache: perto e rápido, mas pode estar velho}
  \begin{columns}[T]
    \begin{column}{0.5\textwidth}
      \begin{definicao}[Cache]
        Cópia de dados guardada perto de quem usa para evitar buscá-los de novo. \textbf{Acerto}: o dado estava lá. \textbf{Falta}: é preciso ir à fonte.
      \end{definicao}
      \scriptsize
      \begin{tabularx}{\linewidth}{@{}l L@{}}
        \toprule
        \cabecalho{Estratégia} & \cabecalho{Como funciona} \\
        \midrule
        Cache primeiro & usa o cache; só vai à rede se não houver \\
        Rede primeiro & tenta a rede; usa o cache se falhar \\
        Mostra e revalida & mostra o cache já e atualiza em seguida (\emph{stale-while-revalidate}) \\
        Só rede & dados que não podem estar velhos (saldo, reserva) \\
        \bottomrule
      \end{tabularx}
    \end{column}
    \begin{column}{0.48\textwidth}
      \begin{conceito}[Invalidação]
        Decidir quando o cache deixou de valer: por tempo (TTL), por evento (o servidor avisou) ou por versão (\emph{ETag}, Módulo 11).
      \end{conceito}
      \begin{pordentro}[Quem sai quando enche]
        A política mais comum é a LRU (\emph{least recently used}): descarta o item usado há mais tempo. Imagens em memória costumam usar esse critério.
      \end{pordentro}
      \begin{historia}
        Phil Karlton brincava que só há duas coisas difíceis na computação: invalidar cache e dar nomes às coisas.
      \end{historia}
    \end{column}
  \end{columns}
\end{frame}

\begin{frame}{\emph{Offline-first}: o banco local como fonte da verdade}
  \begin{columns}[c]
    \begin{column}{0.5\textwidth}
      \centering
      \begin{tikzpicture}[every node/.style={font=\tiny}]
        \node[caixa,font=\tiny,minimum width=2.4cm] (ui) at (0,0) {interface};
        \node[caixaescura,font=\tiny,minimum width=2.4cm] (db) at (0,-1.4) {banco local};
        \node[caixaambar,font=\tiny,minimum width=1.9cm] (fila) at (3.7,-1.4) {fila de envio};
        \node[caixapetroleo,font=\tiny,minimum width=2.4cm] (sync) at (1.6,-2.8) {sincronizador};
        \node[caixa,font=\tiny,minimum width=1.6cm] (api) at (5.0,-2.8) {servidor};
        \draw[seta] ([xshift=-5mm]db.north) -- node[rotulo,left]{observa} ([xshift=-5mm]ui.south);
        \draw[setaambar] ([xshift=5mm]ui.south) -- node[rotulo,right]{grava} ([xshift=5mm]db.north);
        \draw[setaambar] (db) -- node[rotulo,above]{mesma transação} (fila);
        \draw[seta] (fila) -- (sync);
        \draw[seta] ([yshift=1mm]sync.east) -- node[rotulo,above]{envia} ([yshift=1mm]api.west);
        \draw[seta] ([yshift=-1mm]api.west) -- node[rotulo,below]{novidades} ([yshift=-1mm]sync.east);
        \draw[seta] (sync) -- node[rotulo,left]{atualiza} (db);
\end{tikzpicture}
    \end{column}
    \begin{column}{0.48\textwidth}
      \begin{definicao}[\emph{Offline-first}]
        A interface lê e grava \textbf{sempre} no banco local. Um sincronizador troca dados com o servidor quando há rede. Sem rede, o app continua funcionando.
      \end{definicao}
      \begin{conceito}[Fila de envio (\emph{outbox})]
        Cada alteração local vira um registro pendente. O sincronizador envia em ordem, com novas tentativas e espera crescente, mesmo com o app fechado (WorkManager, \cod{workmanager}).
      \end{conceito}
      \begin{atencao}
        Reenvios podem duplicar operações: mande uma chave de idempotência para o servidor reconhecer repetições (Módulo 17).
      \end{atencao}
    \end{column}
  \end{columns}
\end{frame}

\begin{frame}{Sincronização e conflitos}
  \begin{columns}[T]
    \begin{column}{0.52\textwidth}
      \begin{exemplo}[O conflito]
        Ana e Bruno, sem internet, reservam a bancada B3 para as 14\,h. Quando a rede volta, as duas reservas chegam ao servidor. Quem fica com a bancada?
      \end{exemplo}
      \scriptsize
      \begin{tabularx}{\linewidth}{@{}l L@{}}
        \toprule
        \cabecalho{Estratégia} & \cabecalho{Quando usar} \\
        \midrule
        Última gravação vence & dados pessoais pouco disputados; depende de relógios confiáveis \\
        Regra no servidor & recursos disputados: a segunda reserva é recusada \\
        Mesclagem por campo & edições em campos diferentes do mesmo registro \\
        Perguntar ao usuário & textos editados dos dois lados \\
        CRDT & dados colaborativos que se mesclam sozinhos \\
        \bottomrule
      \end{tabularx}
    \end{column}
    \begin{column}{0.46\textwidth}
      \begin{definicao}[Concorrência otimista]
        Cada registro tem uma versão. Quem grava informa a versão que leu; se outra pessoa gravou antes, o servidor recusa e o app resolve o conflito.
      \end{definicao}
      \begin{dica}[Na interface]
        Mostre o que está pendente de envio e avise com clareza quando uma operação for recusada. Ações que não podem esperar (pagar, reservar) devem exigir rede.
      \end{dica}
    \end{column}
  \end{columns}
\end{frame}

\begin{frame}{Dados em repouso: criptografia e segredos}
  \begin{columns}[T]
    \begin{column}{0.5\textwidth}
      \begin{definicao}[Criptografia em repouso]
        Os dados gravados ficam cifrados e só são decifrados com uma chave. Quem copia o armazenamento sem a chave vê apenas bytes sem sentido.
      \end{definicao}
      \scriptsize
      \begin{tabularx}{\linewidth}{@{}l L@{}}
        \toprule
        \cabecalho{Camada} & \cabecalho{O que protege} \\
        \midrule
        Sistema & todo o armazenamento: criptografia por arquivo (Android) e \emph{Data Protection} (iOS) \\
        Cofre de chaves & chaves e segredos em hardware: Keystore (Android), Keychain e Secure Enclave (iOS) \\
        App & banco cifrado (SQLCipher) e campos sensíveis cifrados \\
        \bottomrule
      \end{tabularx}
    \end{column}
    \begin{column}{0.48\textwidth}
      \begin{pordentro}[Por que o hardware importa]
        Chaves criadas no Keystore ou no Secure Enclave não saem do chip: o app pede ``cifre isto'' e recebe o resultado, sem nunca ver a chave.
      \end{pordentro}
      \begin{atencao}[O que nunca fazer]
        Guardar senhas; guardar \emph{tokens} em preferências comuns; deixar chaves de API fixas no código, que pode ser extraído do pacote publicado.
      \end{atencao}
      \begin{dica}
        Pense no risco: aparelho perdido, aparelho com \emph{root} e cópias de segurança na nuvem.
      \end{dica}
    \end{column}
  \end{columns}
\end{frame}

%  (modelagem relacional, SQL, transações e índices estão no Módulo 09)
% =============================================================================
\section{Fundamentos: persistência, cache e sincronização}

\begin{frame}{Memória e armazenamento: velocidade \emph{versus} permanência}
  \begin{columns}[c]
    \begin{column}{0.5\textwidth}
      \centering
      \begin{tikzpicture}[every node/.style={font=\tiny}]
        \foreach \y/\w/\texto/\cor in {3/3.0/{registradores e cache: $\sim$1\,ns}/azulpetroleo,2.25/3.6/{memória RAM: $\sim$100\,ns}/azulpetroleo!80,1.5/4.2/{\emph{flash}: dezenas de µs}/ambar!85,0.75/4.8/{rede móvel: dezenas a centenas de ms}/cinzamedio!70}{
          \fill[\cor,rounded corners=2pt] (-\w/2,\y) rectangle (\w/2,\y+0.6);
          \node[text=white,font=\tiny\bfseries] at (0,\y+0.3) {\texto};
        }
        \draw[seta] (-2.75,3.55) -- node[rotulo,left,text width=1cm,align=right]{mais rápido e menor} (-2.75,0.8);
        \node[rotulo] at (0,0.45) {RAM: volátil \quad \emph{flash}: persistente};
\end{tikzpicture}
    \end{column}
    \begin{column}{0.48\textwidth}
      \begin{definicao}[Persistência]
        Guardar dados de forma que sobrevivam ao fim do processo, ao desligar do aparelho e à falta de rede. No celular, isso significa gravar na memória \emph{flash}.
      \end{definicao}
      \begin{conceito}[Por que persistir em um app]
        O sistema encerra apps em segundo plano (Módulo 07), a rede falha e o usuário espera abrir o app e ver seus dados na hora.
      \end{conceito}
      \begin{pordentro}[A memória \emph{flash}]
        Grava em blocos e se desgasta com o uso. Muitas gravações pequenas custam mais que uma grande: agrupe operações em transações (Módulo 09).
      \end{pordentro}
    \end{column}
  \end{columns}
\end{frame}

\begin{frame}{Onde o sistema guarda os arquivos do app}
  \scriptsize
  \begin{tabularx}{\textwidth}{@{}p{2.5cm} p{2.6cm} p{3.4cm} >{\tiny}L@{}}
    \toprule
    \cabecalho{Finalidade} & \cabecalho{Android} & \cabecalho{iOS} & \scriptsize\cabecalho{Flutter (\cod{path\_provider})} \\
    \midrule
    Dados internos do app & \cod{filesDir}, \cod{databases/} & \cod{Library/Application Support} & \cod{getApplicationSupportDirectory()} \\
    Documentos do usuário & \cod{filesDir} ou MediaStore & \cod{Documents} & \cod{getApplicationDocumentsDirectory()} \\
    Cache descartável & \cod{cacheDir} & \cod{Library/Caches} & \cod{getApplicationCacheDirectory()} \\
    Temporários & \cod{cacheDir} & \cod{tmp} & \cod{getTemporaryDirectory()} \\
    \bottomrule
  \end{tabularx}
  \normalsize
  \vspace{3pt}
  \begin{columns}[T]
    \begin{column}{0.32\textwidth}
      \begin{definicao}[\emph{Sandbox} de arquivos]
        Cada app tem pastas privadas que outros apps não leem. Desinstalar o app apaga tudo.
      \end{definicao}
    \end{column}
    \begin{column}{0.32\textwidth}
      \begin{pordentro}[Cópia de segurança]
        O Android copia até 25\,MB por app para a conta Google; o iOS inclui \cod{Documents} e \cod{Library}, exceto \cod{Caches}.
      \end{pordentro}
    \end{column}
    \begin{column}{0.32\textwidth}
      \begin{atencao}[Cache some]
        O sistema apaga caches quando falta espaço. Nunca guarde ali o que não puder baixar de novo.
      \end{atencao}
    \end{column}
  \end{columns}
\end{frame}

\begin{frame}{Do objeto à tabela}
  \begin{columns}[c]
    \begin{column}{0.46\textwidth}
      \centering
      \begin{tikzpicture}[every node/.style={font=\tiny}]
        \node[caixaescura,font=\tiny,align=left,text width=2.2cm] (o) at (0,0) {\textbf{objeto} \cod{Reserva}\\ \cod{bancada: Bancada}\\ \cod{leituras: List<...>}};
        \node[caixa,font=\tiny,minimum width=1.7cm] (t1) at (3.3,0.9) {tabela \cod{reserva}};
        \node[caixa,font=\tiny,minimum width=1.7cm] (t2) at (3.3,0) {tabela \cod{bancada}};
        \node[caixa,font=\tiny,minimum width=1.7cm] (t3) at (3.3,-0.9) {tabela \cod{leitura}};
        \draw[setaambar] (o.east) -- (t1.west); \draw[setaambar] (o.east) -- (t2.west); \draw[setaambar] (o.east) -- (t3.west);
        \node[rotulo,text width=4.6cm] at (1.6,-1.65) {referências viram chaves estrangeiras; listas viram outra tabela};
      \end{tikzpicture}
      \begin{definicao}[Descompasso objeto-relacional]
        Objetos têm referências, listas e herança; tabelas têm linhas, colunas e chaves. Traduzir de um mundo para o outro é o trabalho do mapeamento.
      \end{definicao}
    \end{column}
    \begin{column}{0.52\textwidth}
      \scriptsize
      \begin{tabularx}{\linewidth}{@{}l L@{}}
        \toprule
        \cabecalho{Abordagem} & \cabecalho{Como é} \\
        \midrule
        SQL à mão & você escreve o SQL e converte as linhas (\cod{sqflite}) \\
        SQL verificado & você escreve o SQL; a ferramenta confere na compilação e gera o código (Room, Drift) \\
        ORM completo & a biblioteca gera o SQL a partir dos objetos \\
        Banco de objetos & guarda objetos sem tabelas (ObjectBox, Realm) \\
        \bottomrule
      \end{tabularx}
      \normalsize
      \begin{dica}
        Room e Drift ficam no meio-termo: SQL explícito, mas com tipos e erros apontados antes de o app rodar.
      \end{dica}
      \begin{atencao}[O problema N+1]
        Buscar uma lista e depois uma consulta por item gera N+1 idas ao banco. Use \cod{JOIN} ou relações carregadas de uma vez.
      \end{atencao}
    \end{column}
  \end{columns}
\end{frame}

\begin{frame}{O esquema também tem versões}
  \centering
  \begin{tikzpicture}[node distance=16mm,every node/.style={font=\scriptsize}]
    \node[caixa,minimum width=1.4cm] (v1) {v1};
    \node[caixa,minimum width=1.4cm,right=of v1] (v2) {v2};
    \node[caixa,minimum width=1.4cm,right=of v2] (v3) {v3};
    \node[caixaescura,minimum width=1.4cm,right=of v3] (v4) {v4 (atual)};
    \draw[seta] (v1) -- node[rotulo,above]{+ coluna} (v2);
    \draw[seta] (v2) -- node[rotulo,above]{nova tabela} (v3);
    \draw[seta] (v3) -- node[rotulo,above]{renomeia} (v4);
    \draw[setaambar,rounded corners] (v1.south) -- ++(0,-4mm) -| node[rotulo,pos=0.25,below]{usuário que pulou versões: o app aplica 1$\rightarrow$2$\rightarrow$3$\rightarrow$4 em sequência} (v4.south);
  \end{tikzpicture}
  \raggedright
  \begin{columns}[T]
    \begin{column}{0.49\textwidth}
      \begin{definicao}[Migração]
        Programa que leva o banco de uma versão do esquema para a seguinte \textbf{sem perder os dados} do usuário. Cada versão publicada do app pode encontrar qualquer versão anterior do banco.
      \end{definicao}
    \end{column}
    \begin{column}{0.49\textwidth}
      \begin{conceito}[Regras de ouro]
        Migrações só andam para a frente e nunca são editadas depois de publicadas; teste a partir de \textbf{cada} versão antiga; apagar e recriar o banco só serve para dados de cache.
      \end{conceito}
    \end{column}
  \end{columns}
\end{frame}

\begin{frame}{Cache: perto e rápido, mas pode estar velho}
  \begin{columns}[T]
    \begin{column}{0.5\textwidth}
      \begin{definicao}[Cache]
        Cópia de dados guardada perto de quem usa para evitar buscá-los de novo. \textbf{Acerto}: o dado estava lá. \textbf{Falta}: é preciso ir à fonte.
      \end{definicao}
      \scriptsize
      \begin{tabularx}{\linewidth}{@{}l L@{}}
        \toprule
        \cabecalho{Estratégia} & \cabecalho{Como funciona} \\
        \midrule
        Cache primeiro & usa o cache; só vai à rede se não houver \\
        Rede primeiro & tenta a rede; usa o cache se falhar \\
        Mostra e revalida & mostra o cache já e atualiza em seguida (\emph{stale-while-revalidate}) \\
        Só rede & dados que não podem estar velhos (saldo, reserva) \\
        \bottomrule
      \end{tabularx}
    \end{column}
    \begin{column}{0.48\textwidth}
      \begin{conceito}[Invalidação]
        Decidir quando o cache deixou de valer: por tempo (TTL), por evento (o servidor avisou) ou por versão (\emph{ETag}, Módulo 11).
      \end{conceito}
      \begin{pordentro}[Quem sai quando enche]
        A política mais comum é a LRU (\emph{least recently used}): descarta o item usado há mais tempo. Imagens em memória costumam usar esse critério.
      \end{pordentro}
      \begin{historia}
        Phil Karlton brincava que só há duas coisas difíceis na computação: invalidar cache e dar nomes às coisas.
      \end{historia}
    \end{column}
  \end{columns}
\end{frame}

\begin{frame}{\emph{Offline-first}: o banco local como fonte da verdade}
  \begin{columns}[c]
    \begin{column}{0.5\textwidth}
      \centering
      \begin{tikzpicture}[every node/.style={font=\tiny}]
        \node[caixa,font=\tiny,minimum width=2.4cm] (ui) at (0,0) {interface};
        \node[caixaescura,font=\tiny,minimum width=2.4cm] (db) at (0,-1.4) {banco local};
        \node[caixaambar,font=\tiny,minimum width=1.9cm] (fila) at (3.7,-1.4) {fila de envio};
        \node[caixapetroleo,font=\tiny,minimum width=2.4cm] (sync) at (1.6,-2.8) {sincronizador};
        \node[caixa,font=\tiny,minimum width=1.6cm] (api) at (5.0,-2.8) {servidor};
        \draw[seta] ([xshift=-5mm]db.north) -- node[rotulo,left]{observa} ([xshift=-5mm]ui.south);
        \draw[setaambar] ([xshift=5mm]ui.south) -- node[rotulo,right]{grava} ([xshift=5mm]db.north);
        \draw[setaambar] (db) -- node[rotulo,above]{mesma transação} (fila);
        \draw[seta] (fila) -- (sync);
        \draw[seta] ([yshift=1mm]sync.east) -- node[rotulo,above]{envia} ([yshift=1mm]api.west);
        \draw[seta] ([yshift=-1mm]api.west) -- node[rotulo,below]{novidades} ([yshift=-1mm]sync.east);
        \draw[seta] (sync) -- node[rotulo,left]{atualiza} (db);
\end{tikzpicture}
    \end{column}
    \begin{column}{0.48\textwidth}
      \begin{definicao}[\emph{Offline-first}]
        A interface lê e grava \textbf{sempre} no banco local. Um sincronizador troca dados com o servidor quando há rede. Sem rede, o app continua funcionando.
      \end{definicao}
      \begin{conceito}[Fila de envio (\emph{outbox})]
        Cada alteração local vira um registro pendente. O sincronizador envia em ordem, com novas tentativas e espera crescente, mesmo com o app fechado (WorkManager, \cod{workmanager}).
      \end{conceito}
      \begin{atencao}
        Reenvios podem duplicar operações: mande uma chave de idempotência para o servidor reconhecer repetições (Módulo 17).
      \end{atencao}
    \end{column}
  \end{columns}
\end{frame}

\begin{frame}{Sincronização e conflitos}
  \begin{columns}[T]
    \begin{column}{0.52\textwidth}
      \begin{exemplo}[O conflito]
        Ana e Bruno, sem internet, reservam a bancada B3 para as 14\,h. Quando a rede volta, as duas reservas chegam ao servidor. Quem fica com a bancada?
      \end{exemplo}
      \scriptsize
      \begin{tabularx}{\linewidth}{@{}l L@{}}
        \toprule
        \cabecalho{Estratégia} & \cabecalho{Quando usar} \\
        \midrule
        Última gravação vence & dados pessoais pouco disputados; depende de relógios confiáveis \\
        Regra no servidor & recursos disputados: a segunda reserva é recusada \\
        Mesclagem por campo & edições em campos diferentes do mesmo registro \\
        Perguntar ao usuário & textos editados dos dois lados \\
        CRDT & dados colaborativos que se mesclam sozinhos \\
        \bottomrule
      \end{tabularx}
    \end{column}
    \begin{column}{0.46\textwidth}
      \begin{definicao}[Concorrência otimista]
        Cada registro tem uma versão. Quem grava informa a versão que leu; se outra pessoa gravou antes, o servidor recusa e o app resolve o conflito.
      \end{definicao}
      \begin{dica}[Na interface]
        Mostre o que está pendente de envio e avise com clareza quando uma operação for recusada. Ações que não podem esperar (pagar, reservar) devem exigir rede.
      \end{dica}
    \end{column}
  \end{columns}
\end{frame}

\begin{frame}{Dados em repouso: criptografia e segredos}
  \begin{columns}[T]
    \begin{column}{0.5\textwidth}
      \begin{definicao}[Criptografia em repouso]
        Os dados gravados ficam cifrados e só são decifrados com uma chave. Quem copia o armazenamento sem a chave vê apenas bytes sem sentido.
      \end{definicao}
      \scriptsize
      \begin{tabularx}{\linewidth}{@{}l L@{}}
        \toprule
        \cabecalho{Camada} & \cabecalho{O que protege} \\
        \midrule
        Sistema & todo o armazenamento: criptografia por arquivo (Android) e \emph{Data Protection} (iOS) \\
        Cofre de chaves & chaves e segredos em hardware: Keystore (Android), Keychain e Secure Enclave (iOS) \\
        App & banco cifrado (SQLCipher) e campos sensíveis cifrados \\
        \bottomrule
      \end{tabularx}
    \end{column}
    \begin{column}{0.48\textwidth}
      \begin{pordentro}[Por que o hardware importa]
        Chaves criadas no Keystore ou no Secure Enclave não saem do chip: o app pede ``cifre isto'' e recebe o resultado, sem nunca ver a chave.
      \end{pordentro}
      \begin{atencao}[O que nunca fazer]
        Guardar senhas; guardar \emph{tokens} em preferências comuns; deixar chaves de API fixas no código, que pode ser extraído do pacote publicado.
      \end{atencao}
      \begin{dica}
        Pense no risco: aparelho perdido, aparelho com \emph{root} e cópias de segurança na nuvem.
      \end{dica}
    \end{column}
  \end{columns}
\end{frame}

%  (modelagem relacional, SQL, transações e índices estão no Módulo 09)
% =============================================================================
\section{Fundamentos: persistência, cache e sincronização}

\begin{frame}{Memória e armazenamento: velocidade \emph{versus} permanência}
  \begin{columns}[c]
    \begin{column}{0.5\textwidth}
      \centering
      \begin{tikzpicture}[every node/.style={font=\tiny}]
        \foreach \y/\w/\texto/\cor in {3/3.0/{registradores e cache: $\sim$1\,ns}/azulpetroleo,2.25/3.6/{memória RAM: $\sim$100\,ns}/azulpetroleo!80,1.5/4.2/{\emph{flash}: dezenas de µs}/ambar!85,0.75/4.8/{rede móvel: dezenas a centenas de ms}/cinzamedio!70}{
          \fill[\cor,rounded corners=2pt] (-\w/2,\y) rectangle (\w/2,\y+0.6);
          \node[text=white,font=\tiny\bfseries] at (0,\y+0.3) {\texto};
        }
        \draw[seta] (-2.75,3.55) -- node[rotulo,left,text width=1cm,align=right]{mais rápido e menor} (-2.75,0.8);
        \node[rotulo] at (0,0.45) {RAM: volátil \quad \emph{flash}: persistente};
\end{tikzpicture}
    \end{column}
    \begin{column}{0.48\textwidth}
      \begin{definicao}[Persistência]
        Guardar dados de forma que sobrevivam ao fim do processo, ao desligar do aparelho e à falta de rede. No celular, isso significa gravar na memória \emph{flash}.
      \end{definicao}
      \begin{conceito}[Por que persistir em um app]
        O sistema encerra apps em segundo plano (Módulo 07), a rede falha e o usuário espera abrir o app e ver seus dados na hora.
      \end{conceito}
      \begin{pordentro}[A memória \emph{flash}]
        Grava em blocos e se desgasta com o uso. Muitas gravações pequenas custam mais que uma grande: agrupe operações em transações (Módulo 09).
      \end{pordentro}
    \end{column}
  \end{columns}
\end{frame}

\begin{frame}{Onde o sistema guarda os arquivos do app}
  \scriptsize
  \begin{tabularx}{\textwidth}{@{}p{2.5cm} p{2.6cm} p{3.4cm} >{\tiny}L@{}}
    \toprule
    \cabecalho{Finalidade} & \cabecalho{Android} & \cabecalho{iOS} & \scriptsize\cabecalho{Flutter (\cod{path\_provider})} \\
    \midrule
    Dados internos do app & \cod{filesDir}, \cod{databases/} & \cod{Library/Application Support} & \cod{getApplicationSupportDirectory()} \\
    Documentos do usuário & \cod{filesDir} ou MediaStore & \cod{Documents} & \cod{getApplicationDocumentsDirectory()} \\
    Cache descartável & \cod{cacheDir} & \cod{Library/Caches} & \cod{getApplicationCacheDirectory()} \\
    Temporários & \cod{cacheDir} & \cod{tmp} & \cod{getTemporaryDirectory()} \\
    \bottomrule
  \end{tabularx}
  \normalsize
  \vspace{3pt}
  \begin{columns}[T]
    \begin{column}{0.32\textwidth}
      \begin{definicao}[\emph{Sandbox} de arquivos]
        Cada app tem pastas privadas que outros apps não leem. Desinstalar o app apaga tudo.
      \end{definicao}
    \end{column}
    \begin{column}{0.32\textwidth}
      \begin{pordentro}[Cópia de segurança]
        O Android copia até 25\,MB por app para a conta Google; o iOS inclui \cod{Documents} e \cod{Library}, exceto \cod{Caches}.
      \end{pordentro}
    \end{column}
    \begin{column}{0.32\textwidth}
      \begin{atencao}[Cache some]
        O sistema apaga caches quando falta espaço. Nunca guarde ali o que não puder baixar de novo.
      \end{atencao}
    \end{column}
  \end{columns}
\end{frame}

\begin{frame}{Do objeto à tabela}
  \begin{columns}[c]
    \begin{column}{0.46\textwidth}
      \centering
      \begin{tikzpicture}[every node/.style={font=\tiny}]
        \node[caixaescura,font=\tiny,align=left,text width=2.2cm] (o) at (0,0) {\textbf{objeto} \cod{Reserva}\\ \cod{bancada: Bancada}\\ \cod{leituras: List<...>}};
        \node[caixa,font=\tiny,minimum width=1.7cm] (t1) at (3.3,0.9) {tabela \cod{reserva}};
        \node[caixa,font=\tiny,minimum width=1.7cm] (t2) at (3.3,0) {tabela \cod{bancada}};
        \node[caixa,font=\tiny,minimum width=1.7cm] (t3) at (3.3,-0.9) {tabela \cod{leitura}};
        \draw[setaambar] (o.east) -- (t1.west); \draw[setaambar] (o.east) -- (t2.west); \draw[setaambar] (o.east) -- (t3.west);
        \node[rotulo,text width=4.6cm] at (1.6,-1.65) {referências viram chaves estrangeiras; listas viram outra tabela};
      \end{tikzpicture}
      \begin{definicao}[Descompasso objeto-relacional]
        Objetos têm referências, listas e herança; tabelas têm linhas, colunas e chaves. Traduzir de um mundo para o outro é o trabalho do mapeamento.
      \end{definicao}
    \end{column}
    \begin{column}{0.52\textwidth}
      \scriptsize
      \begin{tabularx}{\linewidth}{@{}l L@{}}
        \toprule
        \cabecalho{Abordagem} & \cabecalho{Como é} \\
        \midrule
        SQL à mão & você escreve o SQL e converte as linhas (\cod{sqflite}) \\
        SQL verificado & você escreve o SQL; a ferramenta confere na compilação e gera o código (Room, Drift) \\
        ORM completo & a biblioteca gera o SQL a partir dos objetos \\
        Banco de objetos & guarda objetos sem tabelas (ObjectBox, Realm) \\
        \bottomrule
      \end{tabularx}
      \normalsize
      \begin{dica}
        Room e Drift ficam no meio-termo: SQL explícito, mas com tipos e erros apontados antes de o app rodar.
      \end{dica}
      \begin{atencao}[O problema N+1]
        Buscar uma lista e depois uma consulta por item gera N+1 idas ao banco. Use \cod{JOIN} ou relações carregadas de uma vez.
      \end{atencao}
    \end{column}
  \end{columns}
\end{frame}

\begin{frame}{O esquema também tem versões}
  \centering
  \begin{tikzpicture}[node distance=16mm,every node/.style={font=\scriptsize}]
    \node[caixa,minimum width=1.4cm] (v1) {v1};
    \node[caixa,minimum width=1.4cm,right=of v1] (v2) {v2};
    \node[caixa,minimum width=1.4cm,right=of v2] (v3) {v3};
    \node[caixaescura,minimum width=1.4cm,right=of v3] (v4) {v4 (atual)};
    \draw[seta] (v1) -- node[rotulo,above]{+ coluna} (v2);
    \draw[seta] (v2) -- node[rotulo,above]{nova tabela} (v3);
    \draw[seta] (v3) -- node[rotulo,above]{renomeia} (v4);
    \draw[setaambar,rounded corners] (v1.south) -- ++(0,-4mm) -| node[rotulo,pos=0.25,below]{usuário que pulou versões: o app aplica 1$\rightarrow$2$\rightarrow$3$\rightarrow$4 em sequência} (v4.south);
  \end{tikzpicture}
  \raggedright
  \begin{columns}[T]
    \begin{column}{0.49\textwidth}
      \begin{definicao}[Migração]
        Programa que leva o banco de uma versão do esquema para a seguinte \textbf{sem perder os dados} do usuário. Cada versão publicada do app pode encontrar qualquer versão anterior do banco.
      \end{definicao}
    \end{column}
    \begin{column}{0.49\textwidth}
      \begin{conceito}[Regras de ouro]
        Migrações só andam para a frente e nunca são editadas depois de publicadas; teste a partir de \textbf{cada} versão antiga; apagar e recriar o banco só serve para dados de cache.
      \end{conceito}
    \end{column}
  \end{columns}
\end{frame}

\begin{frame}{Cache: perto e rápido, mas pode estar velho}
  \begin{columns}[T]
    \begin{column}{0.5\textwidth}
      \begin{definicao}[Cache]
        Cópia de dados guardada perto de quem usa para evitar buscá-los de novo. \textbf{Acerto}: o dado estava lá. \textbf{Falta}: é preciso ir à fonte.
      \end{definicao}
      \scriptsize
      \begin{tabularx}{\linewidth}{@{}l L@{}}
        \toprule
        \cabecalho{Estratégia} & \cabecalho{Como funciona} \\
        \midrule
        Cache primeiro & usa o cache; só vai à rede se não houver \\
        Rede primeiro & tenta a rede; usa o cache se falhar \\
        Mostra e revalida & mostra o cache já e atualiza em seguida (\emph{stale-while-revalidate}) \\
        Só rede & dados que não podem estar velhos (saldo, reserva) \\
        \bottomrule
      \end{tabularx}
    \end{column}
    \begin{column}{0.48\textwidth}
      \begin{conceito}[Invalidação]
        Decidir quando o cache deixou de valer: por tempo (TTL), por evento (o servidor avisou) ou por versão (\emph{ETag}, Módulo 11).
      \end{conceito}
      \begin{pordentro}[Quem sai quando enche]
        A política mais comum é a LRU (\emph{least recently used}): descarta o item usado há mais tempo. Imagens em memória costumam usar esse critério.
      \end{pordentro}
      \begin{historia}
        Phil Karlton brincava que só há duas coisas difíceis na computação: invalidar cache e dar nomes às coisas.
      \end{historia}
    \end{column}
  \end{columns}
\end{frame}

\begin{frame}{\emph{Offline-first}: o banco local como fonte da verdade}
  \begin{columns}[c]
    \begin{column}{0.5\textwidth}
      \centering
      \begin{tikzpicture}[every node/.style={font=\tiny}]
        \node[caixa,font=\tiny,minimum width=2.4cm] (ui) at (0,0) {interface};
        \node[caixaescura,font=\tiny,minimum width=2.4cm] (db) at (0,-1.4) {banco local};
        \node[caixaambar,font=\tiny,minimum width=1.9cm] (fila) at (3.7,-1.4) {fila de envio};
        \node[caixapetroleo,font=\tiny,minimum width=2.4cm] (sync) at (1.6,-2.8) {sincronizador};
        \node[caixa,font=\tiny,minimum width=1.6cm] (api) at (5.0,-2.8) {servidor};
        \draw[seta] ([xshift=-5mm]db.north) -- node[rotulo,left]{observa} ([xshift=-5mm]ui.south);
        \draw[setaambar] ([xshift=5mm]ui.south) -- node[rotulo,right]{grava} ([xshift=5mm]db.north);
        \draw[setaambar] (db) -- node[rotulo,above]{mesma transação} (fila);
        \draw[seta] (fila) -- (sync);
        \draw[seta] ([yshift=1mm]sync.east) -- node[rotulo,above]{envia} ([yshift=1mm]api.west);
        \draw[seta] ([yshift=-1mm]api.west) -- node[rotulo,below]{novidades} ([yshift=-1mm]sync.east);
        \draw[seta] (sync) -- node[rotulo,left]{atualiza} (db);
\end{tikzpicture}
    \end{column}
    \begin{column}{0.48\textwidth}
      \begin{definicao}[\emph{Offline-first}]
        A interface lê e grava \textbf{sempre} no banco local. Um sincronizador troca dados com o servidor quando há rede. Sem rede, o app continua funcionando.
      \end{definicao}
      \begin{conceito}[Fila de envio (\emph{outbox})]
        Cada alteração local vira um registro pendente. O sincronizador envia em ordem, com novas tentativas e espera crescente, mesmo com o app fechado (WorkManager, \cod{workmanager}).
      \end{conceito}
      \begin{atencao}
        Reenvios podem duplicar operações: mande uma chave de idempotência para o servidor reconhecer repetições (Módulo 17).
      \end{atencao}
    \end{column}
  \end{columns}
\end{frame}

\begin{frame}{Sincronização e conflitos}
  \begin{columns}[T]
    \begin{column}{0.52\textwidth}
      \begin{exemplo}[O conflito]
        Ana e Bruno, sem internet, reservam a bancada B3 para as 14\,h. Quando a rede volta, as duas reservas chegam ao servidor. Quem fica com a bancada?
      \end{exemplo}
      \scriptsize
      \begin{tabularx}{\linewidth}{@{}l L@{}}
        \toprule
        \cabecalho{Estratégia} & \cabecalho{Quando usar} \\
        \midrule
        Última gravação vence & dados pessoais pouco disputados; depende de relógios confiáveis \\
        Regra no servidor & recursos disputados: a segunda reserva é recusada \\
        Mesclagem por campo & edições em campos diferentes do mesmo registro \\
        Perguntar ao usuário & textos editados dos dois lados \\
        CRDT & dados colaborativos que se mesclam sozinhos \\
        \bottomrule
      \end{tabularx}
    \end{column}
    \begin{column}{0.46\textwidth}
      \begin{definicao}[Concorrência otimista]
        Cada registro tem uma versão. Quem grava informa a versão que leu; se outra pessoa gravou antes, o servidor recusa e o app resolve o conflito.
      \end{definicao}
      \begin{dica}[Na interface]
        Mostre o que está pendente de envio e avise com clareza quando uma operação for recusada. Ações que não podem esperar (pagar, reservar) devem exigir rede.
      \end{dica}
    \end{column}
  \end{columns}
\end{frame}

\begin{frame}{Dados em repouso: criptografia e segredos}
  \begin{columns}[T]
    \begin{column}{0.5\textwidth}
      \begin{definicao}[Criptografia em repouso]
        Os dados gravados ficam cifrados e só são decifrados com uma chave. Quem copia o armazenamento sem a chave vê apenas bytes sem sentido.
      \end{definicao}
      \scriptsize
      \begin{tabularx}{\linewidth}{@{}l L@{}}
        \toprule
        \cabecalho{Camada} & \cabecalho{O que protege} \\
        \midrule
        Sistema & todo o armazenamento: criptografia por arquivo (Android) e \emph{Data Protection} (iOS) \\
        Cofre de chaves & chaves e segredos em hardware: Keystore (Android), Keychain e Secure Enclave (iOS) \\
        App & banco cifrado (SQLCipher) e campos sensíveis cifrados \\
        \bottomrule
      \end{tabularx}
    \end{column}
    \begin{column}{0.48\textwidth}
      \begin{pordentro}[Por que o hardware importa]
        Chaves criadas no Keystore ou no Secure Enclave não saem do chip: o app pede ``cifre isto'' e recebe o resultado, sem nunca ver a chave.
      \end{pordentro}
      \begin{atencao}[O que nunca fazer]
        Guardar senhas; guardar \emph{tokens} em preferências comuns; deixar chaves de API fixas no código, que pode ser extraído do pacote publicado.
      \end{atencao}
      \begin{dica}
        Pense no risco: aparelho perdido, aparelho com \emph{root} e cópias de segurança na nuvem.
      \end{dica}
    \end{column}
  \end{columns}
\end{frame}


% =============================================================================
\section{Onde guardar cada dado}

\begin{frame}{Opções de armazenamento}
  \footnotesize
  \begin{tabularx}{\textwidth}{@{}p{3.6cm} L L@{}}
    \toprule
    \cabecalho{Opção} & \cabecalho{Para quê} & \cabecalho{Exemplo no LabMonitor} \\
    \midrule
    \cod{shared\_preferences} & pares chave-valor simples & tema escuro, último laboratório aberto \\
    \cod{flutter\_secure\_storage} & segredos cifrados (Keychain/Keystore) & token de acesso, chave da API \\
    Arquivos (\cod{path\_provider}) & documentos, imagens, exportações & relatório CSV, fotos de equipamentos \\
    \cod{sqflite} & SQLite com SQL escrito à mão & tabelas simples, controle total do SQL \\
    \textbf{Drift} & SQLite tipado, reativo, com migrações & histórico de leituras dos sensores \\
    \bottomrule
  \end{tabularx}
  \vspace{4pt}
  \begin{columns}[T]
    \begin{column}{0.47\textwidth}
      \begin{atencao}
        Nunca guarde tokens ou senhas em \cod{shared\_preferences}: o arquivo não é cifrado.
      \end{atencao}
    \end{column}
    \begin{column}{0.47\textwidth}
      \begin{dica}
        Bancos NoSQL locais (ObjectBox, Hive CE) também existem; confira a manutenção do pacote no pub.dev antes de adotar.
      \end{dica}
    \end{column}
  \end{columns}
\end{frame}

% =============================================================================
\section{Chave-valor, segredos e arquivos}

\begin{frame}[fragile]{Preferências com \cod{shared\_preferences}}
  \begin{columns}[T]
    \begin{column}{0.54\textwidth}
\begin{codigo}[title=data/services/preferencias\_service.dart]{Dart}
import 'package:shared_preferences/shared_preferences.dart';

class PreferenciasService {
  final _prefs = SharedPreferencesAsync();

  Future<bool> modoEscuro() async =>
      await _prefs.getBool('modoEscuro') ?? false;

  Future<void> definirModoEscuro(bool v) =>
      _prefs.setBool('modoEscuro', v);

  Future<String?> ultimoLab() => _prefs.getString('ultimoLab');

  Future<void> limpar() => _prefs.clear();
}
\end{codigo}
    \end{column}
    \begin{column}{0.44\textwidth}
      \begin{conceito}[Duas APIs atuais]
        \begin{itemize}
          \item \cod{SharedPreferencesAsync}: toda leitura vai à plataforma; sempre atualizada.
          \item \cod{SharedPreferencesWithCache}: carrega um conjunto de chaves (\cod{allowList}) e lê de forma síncrona.
        \end{itemize}
      \end{conceito}
      \begin{dica}
        A API antiga \cod{SharedPreferences.getInstance()} ainda funciona, mas a documentação recomenda as duas novas.
      \end{dica}
    \end{column}
  \end{columns}
\end{frame}

\begin{frame}[fragile]{Segredos e arquivos}
  \begin{columns}[T]
    \begin{column}{0.5\textwidth}
\begin{codigo}[title=Armazenamento seguro]{Dart}
const _seguro = FlutterSecureStorage();

Future<void> salvarToken(String jwt) =>
    _seguro.write(key: 'token', value: jwt);

Future<String?> lerToken() =>
    _seguro.read(key: 'token');

Future<void> sair() =>
    _seguro.delete(key: 'token');
\end{codigo}
      \begin{conceito}[Por baixo]
        iOS: \textbf{Keychain}. Android: chave no \textbf{Android Keystore} cifrando os dados. Web: sem garantia equivalente.
      \end{conceito}
    \end{column}
    \begin{column}{0.48\textwidth}
\begin{codigo}[title=Arquivos com path\_provider]{Dart}
Future<File> _arquivo() async {
  final dir = await getApplicationDocumentsDirectory();
  return File('${dir.path}/leituras.csv');
}

Future<void> anexar(Leitura l) async {
  final f = await _arquivo();
  await f.writeAsString(
    '${l.em.toIso8601String()};${l.valor}\n',
    mode: FileMode.append);
}
\end{codigo}
      \begin{dica}
        \cod{getTemporaryDirectory()} para cache que o sistema pode apagar; \cod{getApplicationSupportDirectory()} para arquivos internos.
      \end{dica}
    \end{column}
  \end{columns}
\end{frame}

% =============================================================================
\section{SQLite com sqflite}

\begin{frame}[fragile]{Abrindo o banco e criando tabelas}
  \begin{columns}[T]
    \begin{column}{0.56\textwidth}
\begin{codigo}[title=data/services/banco\_service.dart]{Dart}
import 'package:path/path.dart';
import 'package:sqflite/sqflite.dart';

Future<Database> abrirBanco() async => openDatabase(
  join(await getDatabasesPath(), 'labmonitor.db'),
  version: 2,
  onCreate: (db, versao) async {
    await db.execute('''
      CREATE TABLE leitura (
        id     INTEGER PRIMARY KEY AUTOINCREMENT,
        sensor TEXT NOT NULL,
        valor  REAL NOT NULL,
        em     INTEGER NOT NULL)''');
    await db.execute(
      'CREATE INDEX idx_leitura ON leitura(sensor, em)');
  },
  onUpgrade: (db, de, para) async {
    if (de < 2) await db.execute(
      'ALTER TABLE leitura ADD COLUMN unidade TEXT');
  },
);
\end{codigo}
    \end{column}
    \begin{column}{0.42\textwidth}
      \begin{conceito}[Ciclo de versões]
        \cod{onCreate} roda na primeira instalação; \cod{onUpgrade} roda quando a \cod{version} aumenta. Cada mudança de esquema ganha um bloco \cod{if (de < N)}.
      \end{conceito}
      \begin{atencao}
        Em instalações novas o \cod{onCreate} precisa criar o esquema \textbf{mais recente} (inclua aqui a coluna \cod{unidade}).
      \end{atencao}
      \begin{dica}
        Datas como \cod{INTEGER} (milissegundos desde 1970) ordenam e comparam rápido.
      \end{dica}
    \end{column}
  \end{columns}
\end{frame}

\begin{frame}[fragile]{Consultas, transações e segurança}
  \begin{columns}[T]
    \begin{column}{0.56\textwidth}
\begin{codigo}[title=CRUD com sqflite]{Dart}
Future<int> inserir(Database db, Leitura l) =>
    db.insert('leitura', l.paraMapa());

Future<List<Leitura>> ultimas(Database db, String s) async {
  final linhas = await db.query('leitura',
      where: 'sensor = ?', whereArgs: [s],
      orderBy: 'em DESC', limit: 50);
  return linhas.map(Leitura.deMapa).toList();
}

Future<void> importar(Database db, List<Leitura> ls) =>
    db.transaction((txn) async {
      final lote = txn.batch();
      for (final l in ls) lote.insert('leitura', l.paraMapa());
      await lote.commit(noResult: true);
    });
\end{codigo}
    \end{column}
    \begin{column}{0.42\textwidth}
      \begin{erro}[SQL injection]
        \cod{db.rawQuery("... WHERE sensor = '\$s'")} deixa o usuário escrever SQL. Use sempre \cod{?} com \cod{whereArgs} ou \cod{arguments}.
      \end{erro}
      \begin{conceito}[Transação + batch]
        Tudo ou nada (ACID, Módulo 09) e muito mais rápido do que mil \cod{insert} isolados.
      \end{conceito}
    \end{column}
  \end{columns}
\end{frame}

% =============================================================================
\section{Drift: SQLite tipado e reativo}

\begin{frame}[fragile]{Definindo tabelas e o banco}
  \begin{columns}[T]
    \begin{column}{0.58\textwidth}
\begin{codigo}[title=data/services/banco.dart]{Dart}
import 'package:drift/drift.dart';
import 'package:drift_flutter/drift_flutter.dart';
part 'banco.g.dart';

@TableIndex(name: 'idx_leitura', columns: {#sensor, #em})
class Leituras extends Table {
  IntColumn get id => integer().autoIncrement()();
  TextColumn get sensor => text().withLength(max: 40)();
  RealColumn get valor => real()();
  DateTimeColumn get em => dateTime()();
}

@DriftDatabase(tables: [Leituras])
class BancoApp extends _$BancoApp {
  BancoApp([QueryExecutor? e])
      : super(e ?? driftDatabase(name: 'labmonitor'));

  @override
  int get schemaVersion => 1;
}
\end{codigo}
    \end{column}
    \begin{column}{0.4\textwidth}
\begin{terminal}
$ flutter pub add drift drift_flutter
$ flutter pub add -d drift_dev build_runner
$ dart run build_runner build
\end{terminal}
      \begin{conceito}[O que é gerado]
        A classe de dados \cod{Leitura}, o \cod{LeiturasCompanion} para inserções e as consultas tipadas. Erros de nome de coluna viram erros de compilação.
      \end{conceito}
      \begin{dica}
        O construtor opcional permite passar um banco em memória nos testes.
      \end{dica}
    \end{column}
  \end{columns}
\end{frame}

\begin{frame}[fragile]{Consultas e streams}
  \begin{columns}[T]
    \begin{column}{0.58\textwidth}
\begin{codigo}[title=Dentro de BancoApp (ou de um DAO)]{Dart}
Future<int> registrar(String sensor, double v) =>
    into(leituras).insert(LeiturasCompanion.insert(
        sensor: sensor, valor: v, em: DateTime.now()));

Stream<List<Leitura>> observar(String sensor) =>
    (select(leituras)
          ..where((t) => t.sensor.equals(sensor))
          ..orderBy([(t) => OrderingTerm.desc(t.em)])
          ..limit(50))
        .watch();

Future<int> apagarAntigas(DateTime limite) =>
    (delete(leituras)
          ..where((t) => t.em.isSmallerThanValue(limite)))
        .go();

Future<double?> media(String sensor) => customSelect(
    'SELECT AVG(valor) AS m FROM leituras WHERE sensor = ?',
    variables: [Variable(sensor)], readsFrom: {leituras},
  ).map((r) => r.read<double?>('m')).getSingle();
\end{codigo}
    \end{column}
    \begin{column}{0.4\textwidth}
      \begin{conceito}[\cod{.watch()}]
        Devolve um \cod{Stream} que emite de novo sempre que a tabela muda. A tela se atualiza sozinha quando chega uma leitura nova.
      \end{conceito}
      \begin{dica}
        Também é possível escrever as consultas em SQL puro em arquivos \cod{.drift}; o Drift as valida e gera métodos Dart.
      \end{dica}
    \end{column}
  \end{columns}
\end{frame}

\begin{frame}[fragile]{Migrações com Drift}
  \begin{columns}[T]
    \begin{column}{0.56\textwidth}
\begin{codigo}[title=Versão 2: nova coluna]{Dart}
class Leituras extends Table {
  // ... colunas anteriores
  TextColumn get unidade => text().nullable()();
}

@override
int get schemaVersion => 2;

@override
MigrationStrategy get migration => MigrationStrategy(
  onCreate: (m) => m.createAll(),
  onUpgrade: (m, de, para) async {
    if (de < 2) {
      await m.addColumn(leituras, leituras.unidade);
    }
  },
);
\end{codigo}
    \end{column}
    \begin{column}{0.42\textwidth}
\begin{terminal}
$ dart run drift_dev make-migrations
\end{terminal}
      \begin{conceito}[Migrações verificadas]
        O comando salva o esquema de cada versão e gera migrações passo a passo e testes que conferem se a versão 1 migrada é igual à versão 2 criada do zero.
      \end{conceito}
      \begin{atencao}
        Uma migração errada pode apagar dados de quem já usa o app. Teste sempre a atualização, não só a instalação nova.
      \end{atencao}
    \end{column}
  \end{columns}
\end{frame}

% =============================================================================
\section{Banco local na arquitetura}

\begin{frame}[fragile]{Repositório \emph{offline-first}}
  \begin{columns}[T]
    \begin{column}{0.56\textwidth}
\begin{codigo}[title=data/repositories/leitura\_repository.dart]{Dart}
class LeituraRepository {
  LeituraRepository({required this.banco, required this.api});
  final BancoApp banco;
  final ApiService api;

  // a UI sempre lê do banco local
  Stream<List<Leitura>> observar(String sensor) =>
      banco.observar(sensor);

  // a rede só atualiza o banco
  Future<Result<void>> sincronizar(String sensor) async {
    try {
      final novas = await api.leituras(sensor);
      await banco.batch((b) => b.insertAll(banco.leituras,
          novas.map((d) => d.paraCompanion())));
      return const Ok(null);
    } on Exception catch (e) { return Falha(e); }
  }
}
\end{codigo}
    \end{column}
    \begin{column}{0.42\textwidth}
      \centering
      \begin{tikzpicture}[node distance=5mm]
        \node[caixa] (ui) {View model};
        \node[caixaescura,below=of ui] (db) {Banco local (Drift)};
        \node[caixa,below=of db] (api) {API remota};
        \draw[setaambar] (db) -- node[rotulo,right]{\tiny stream} (ui);
        \draw[seta] (api) -- node[rotulo,right]{\tiny sincroniza} (db);
      \end{tikzpicture}
      \begin{conceito}[Fonte única da verdade]
        A tela observa o banco; a rede apenas o atualiza. Sem internet, o app continua mostrando os últimos dados.
      \end{conceito}
    \end{column}
  \end{columns}
\end{frame}

\begin{frame}[fragile]{Testando o banco}
  \begin{columns}[T]
    \begin{column}{0.56\textwidth}
\begin{codigo}[title=test/banco\_test.dart]{Dart}
import 'package:drift/native.dart';
import 'package:flutter_test/flutter_test.dart';

void main() {
  late BancoApp banco;
  setUp(() => banco = BancoApp(NativeDatabase.memory()));
  tearDown(() => banco.close());

  test('observar emite a leitura registrada', () async {
    await banco.registrar('temp', 23.5);
    final lista = await banco.observar('temp').first;
    expect(lista.single.valor, 23.5);
  });
}
\end{codigo}
    \end{column}
    \begin{column}{0.42\textwidth}
      \begin{conceito}[Banco em memória]
        \cod{NativeDatabase.memory()} cria um SQLite real, descartado a cada teste. Rápido e sem efeitos colaterais.
      \end{conceito}
      \begin{atencao}[Web]
        \cod{sqflite} não roda na web. O Drift funciona na web com WebAssembly, mas exige configuração extra.
      \end{atencao}
    \end{column}
  \end{columns}
\end{frame}

\begin{frame}{sqflite ou Drift?}
  \footnotesize
  \begin{tabularx}{\textwidth}{@{}l L L@{}}
    \toprule
    & \cabecalho{sqflite} & \cabecalho{Drift} \\
    \midrule
    Consultas & SQL em strings & API Dart tipada ou SQL validado \\
    Erros de esquema & em tempo de execução & em tempo de compilação \\
    Reatividade & manual & \cod{.watch()} embutido \\
    Migrações & \cod{onUpgrade} manual & \cod{MigrationStrategy} + verificação gerada \\
    Geração de código & não & sim (\cod{build\_runner}) \\
    Equivalente Kotlin & SQLite direto & Room \\
    \bottomrule
  \end{tabularx}
  \vspace{4pt}
  \begin{caixaambar}
    Para aprender SQL, comece com \dest{sqflite}. Para apps reais com várias tabelas e telas reativas, prefira \dest{Drift}.
  \end{caixaambar}
\end{frame}

% =============================================================================
\section{Fechamento}

\begin{frame}{Resumo}
  \begin{resumo}
    \begin{itemize}
      \item Preferências em \cod{shared\_preferences}; segredos em \cod{flutter\_secure\_storage}; documentos em arquivos.
      \item \cod{sqflite} dá acesso direto ao SQLite: use \cod{?}, transações e \cod{batch}.
      \item Drift gera código tipado, oferece consultas reativas e migrações verificadas.
      \item O banco local é a fonte da verdade; a rede apenas o sincroniza.
      \item Teste com banco em memória e teste também as migrações.
    \end{itemize}
  \end{resumo}
\end{frame}

\begin{frame}{Glossário (1/2)}
  \begin{glossario}[3.8cm]
    Persistência & guardar dados que sobrevivem ao fim do processo e ao desligar \\
    \emph{Sandbox} de arquivos & pastas privadas do app, apagadas na desinstalação \\
    \cod{shared\_preferences} & pares chave-valor simples, sem criptografia \\
    \cod{flutter\_secure\_storage} & segredos cifrados no Keychain ou com o Keystore \\
    \cod{path\_provider} & localiza as pastas do app em cada plataforma \\
    \cod{sqflite} & acesso ao SQLite com SQL escrito à mão \\
    Drift & SQLite tipado e reativo, com código gerado \\
    Migração & leva o banco de uma versão do esquema à seguinte sem perder dados \\
    Problema N+1 & uma consulta para a lista e mais uma para cada item \\
  \end{glossario}
\end{frame}

\begin{frame}{Glossário (2/2)}
  \begin{glossario}
    Cache & cópia guardada perto de quem usa, para não buscar de novo \\
    Invalidação & decidir quando o cache deixou de valer (TTL, evento, versão) \\
    LRU & política que descarta o item usado há mais tempo \\
    \emph{Offline-first} & a interface lê e grava no banco local; a sincronização vem depois \\
    Fila de envio & alterações locais pendentes de envio ao servidor (\emph{outbox}) \\
    Concorrência otimista & a gravação é recusada se a versão lida já mudou \\
    Keystore e Keychain & cofres de chaves do Android e do iOS, protegidos por hardware \\
  \end{glossario}
\end{frame}

\begin{frame}{Questões de revisão}
  \begin{questoes}
    \item Por que um app precisa persistir dados mesmo quando tem internet?
    \item Onde guardar a preferência de tema, o \emph{token} de acesso, o histórico de leituras e as miniaturas baixadas?
    \item Quando preferir Drift a \cod{sqflite}?
    \item Um usuário atualiza da versão 1 direto para a 4 do app. O que a migração do Drift precisa tratar?
    \item Explique o problema N+1 e como evitá-lo.
    \item Descreva o caminho de uma reserva feita sem internet em uma arquitetura \emph{offline-first}.
    \item Dois alunos reservam a mesma bancada sem internet. Que estratégia de conflito usar e por quê?
    \item Por que não guardar o \emph{token} em \cod{shared\_preferences}?
  \end{questoes}
\end{frame}

\begin{frame}{Exercícios}
  \begin{exercicio}[Exercício 1 --- Preferências]
    Salve o tema escolhido e o último laboratório aberto e restaure-os ao abrir o app.
  \end{exercicio}
  \begin{exercicio}[Exercício 2 --- sqflite]
    Crie as tabelas \cod{laboratorio} e \cod{equipamento} (1:N) e liste os equipamentos com um \cod{JOIN}.
  \end{exercicio}
  \begin{exercicio}[Exercício 3 --- Migração]
    Adicione a coluna \cod{unidade} na versão 2 do banco Drift e escreva um teste que migra da versão 1.
  \end{exercicio}
  \begin{projeto}[LabMonitor --- Etapa 10]
    Histórico de leituras em Drift observado pelo painel, preferências com \cod{SharedPreferencesAsync} e token em armazenamento seguro.
  \end{projeto}
\end{frame}

\begin{frame}{Para aprofundar}
  \small
  \begin{itemize}
    \item Persistência no Flutter: \url{https://docs.flutter.dev/cookbook/persistence}
    \item shared\_preferences: \url{https://pub.dev/packages/shared_preferences}
    \item flutter\_secure\_storage: \url{https://pub.dev/packages/flutter_secure_storage}
    \item sqflite: \url{https://pub.dev/packages/sqflite}
    \item Drift: \url{https://drift.simonbinder.eu}
  \end{itemize}
\end{frame}

\begin{frame}{Respostas das questões de revisão}
  \begin{respostas}
    \item O sistema encerra apps em segundo plano, a rede falha e o usuário quer ver os dados na hora: o banco local serve de cache e de fonte da verdade.
    \item Tema: \cod{shared\_preferences}. \emph{Token}: \cod{flutter\_secure\_storage}. Leituras: Drift ou \cod{sqflite}. Miniaturas: \cod{getApplicationCacheDirectory()}.
    \item Quando se quer consultas tipadas e verificadas na geração de código, \emph{streams} reativas e migrações assistidas; \cod{sqflite} atende bancos pequenos com controle total do SQL.
    \item O \cod{onUpgrade} deve aplicar todos os passos de \cod{de} até \cod{para} em sequência (2, 3 e 4), e cada passo precisa de teste a partir da versão antiga (\cod{make-migrations}).
    \item Buscar a lista e depois fazer uma consulta por item gera N+1 idas ao banco. Use um \cod{JOIN} ou carregue os itens relacionados de uma vez.
    \item A reserva é gravada no banco como pendente, junto com um registro na fila de envio, na mesma transação. A tela já a mostra; uma tarefa em segundo plano a envia quando houver rede e atualiza o estado.
    \item Regra no servidor: a bancada é um recurso disputado. A primeira reserva confirmada vence, a segunda é recusada e o app avisa o usuário.
    \item O arquivo não é cifrado pelo app e pode ir para cópias de segurança ou ser lido em aparelhos com \emph{root}. Segredos vão para o Keychain ou para o Keystore, via \cod{flutter\_secure\_storage}.
  \end{respostas}
\end{frame}

\encerramento{Módulo 11 --- Rede e APIs: HTTP, JSON, \cod{http} e Dio.}

\end{document}
